\documentclass[11pt]{article}
\usepackage[margin=1in,letterpaper]{geometry}
\usepackage{amsmath,amssymb,booktabs,array,xcolor}
\usepackage{textcomp}
\usepackage[hidelinks]{hyperref}
\usepackage{graphicx}
\graphicspath{{paper_figures/}{./}{../paper_figures/}}
\usepackage{tabularx}
\usepackage{caption}
\captionsetup{font=small,labelfont=bf,skip=4pt}
\usepackage{placeins}
\renewcommand{\topfraction}{0.9}
\renewcommand{\bottomfraction}{0.85}
\renewcommand{\textfraction}{0.07}
\renewcommand{\floatpagefraction}{0.75}
\setcounter{topnumber}{4}
\setcounter{bottomnumber}{4}
\setcounter{totalnumber}{8}
\emergencystretch=2.5em

\newcommand{\lamsq}{\lambda^{2}}
\newcommand{\ba}{\mbox{b/a}}
\newcommand{\winend}{\mbox{\textsection}}

\title{\bfseries Exact Radial-Determinant Frequencies of Annular and
Circular Plates: the Closed-Ring and Solid-Disk Limits of the
Sector Wave-Function Method}
\author{Garrek McCune$^{a,*}$, Daniel Stutts$^{a}$\\
\small\itshape $^a$Mechanical and Aerospace Engineering,\\
Missouri University of Science and Technology, Rolla, MO 65409, USA\\
\small\upshape $^*$Corresponding author. E-mail: ghmkfh@umsystem.edu}
\date{DRAFT --- \today}
% Supplementary Material mapping (PAPER3_RING_DISK_SUPPLEMENTARY.tex):
%   S.1  Bessel span of the four Frobenius branches
%   S.2  disk regularity, the inner jet, and the rank-3 trap
%   S.3  conversion algebra (one flexural number, one in-plane number)
%   S.4  full match lists at the second Poisson ratio, and remaining Vogel
%   S.5  cluster job provenance
%   S.6  Southwell inverse-b/a table and predicted-versus-observed figure
% Main text contains no job numbers.

\begin{document}
\maketitle

\begin{abstract}
The Seok--Tiersten--Scarton wave-function method for annular-sector plates
has two classical limits that never form the two-edge Galerkin determinant.
At sector angle $2\Theta=2\pi$ there are no radial walls: circumferential
order is an integer $n$, and each frequency is a zero of a $4\times4$ radial
edge-determinant already present in the sector assembler. At inner radius
$R_i=0$ the two origin-singular branches drop out and the same operator
becomes a $2\times2$ at the outer radius --- the classical Kirchhoff disk.
Because there are no $\theta$-edges, Vogel and Skinner's nine inner/outer
combinations are different row pairs of one matrix. Frequencies are found
from $\det L(n,\Omega)=0$ from scratch; published tables enter only post
hoc. A frozen sign-flip ladder replaces the sector residual bar, which does
not travel here. The nine Vogel tables are recovered at 704/710
MATCH${}+{}$WEAK under a frozen $1\%/3\%$ bar. The solid-disk $2\times2$
recovers Narita Table~2(b) at 4/4 and SP-160 Table~2.5 at 6/6; a
$4\times4$ at $R_i=0$ has rank three for $n\ge1$ and is not the disk.
Southwell's Table~2.31 scores 13/18; independent SHELL281 finite elements
at $b/a=0.840$, $n=0$ give $\lamsq=132.549$ against the solver's $132.990$
($+0.333\%$) and Southwell's printed $81.0$ ($-63.6\%$). Polar-orthotropic
free--free rings recover Narita Table~3 at 10/10. In-plane free--free rings
recover Irie, Yamada and Muramoto's Table~2 at 48/48. Default clamped-free
sector assembly is unchanged.
\end{abstract}

\noindent\textbf{Keywords:} annular plate; circular plate; free vibration;
exact solution; radial determinant; mixed boundary conditions; polar
orthotropy; in-plane vibration.

\section{Introduction}\label{sec:intro}
Annular and circular plates are the $2\Theta=2\pi$ and $R_i=0$ limits of the
annular-sector geometry treated in the companion paper
\cite{mccune2026annular}. Classical Kirchhoff tables for those limits have
existed for decades: Leissa's monograph \cite{leissaSP160} compiles the exact
annular results of Vogel and Skinner \cite{vogelSkinner}, Southwell's
inner-clamped annulus as compiled in SP-160 Table~2.31 \cite{southwell},
and the axisymmetric work of Raju \cite{raju}; Narita \cite{narita1984} added
five-figure free--free values for both the annulus and the solid disk; Irie,
Yamada and Muramoto \cite{irie1984} tabulated the in-plane annular spectrum.
The contribution here is not another Bessel catalogue. It is the reduction of
a \emph{validated sector solver} onto those tables, in both motions, with a
screen that does not use the sector residual bar, together with all nine
inner/outer edge combinations that the sector corner problem blocks, and a
polar-orthotropic free--free hold-out on the same operator.

Seok, Tiersten and Scarton \cite{seok1,seok2,seok3,seok4} gave
exact-wave-function solutions for annular-sector and rectangular plates,
out-of-plane and in-plane, at clamped-free. The first companion paper
\cite{mccune2026annular} extended the annular-sector flexural assembler to
completely free sectors and built a residual screen against independent finite
elements; the second \cite{mccune2026rect} did the rectangular analogue, where
the Kirchhoff corner jump is load-bearing. Neither takes the $2\pi$ or $R_i=0$
limit as a production entry point. That limit is cheaper than the sector:
circumferential order is an integer, the characteristic matrix is $4\times4$
(ring) or $2\times2$ (disk), there are no $\theta$-walls, and there is no
corner jump. It is also the elastic forward model that a later piezoelectric
identification will need on the free--free and clamped-free ring.

The detector remains paper-faithful. Frequencies are zeros of
$\det L(n,\Omega)=0$ built from the same Frobenius radial content as the
sector assembler. Published values never seed a search; they enter only when a
converged root is scored against them. The sign-flip ladder described in
\S\ref{sec:screen} was frozen before the production sweeps and was not retuned
to admit a published match. Independent cyclic-symmetry SHELL281 calculations
\cite{ansys} are a third method, not a fitted target.

What this paper is not: McGee's solid \emph{sector} (that is the sector
assembler with a coordinate singularity); elliptical plates; simply supported
or guided sector edges (the open corner-jump problem); piezoelectric coupling,
mentioned here only prospectively, as future work; or a splice into either
companion paper.

\section{Reductions}\label{sec:reduce}

\subsection{Closed ring, out-of-plane}\label{sec:ring}
Requiring $2\pi$-periodicity of the transverse displacement forces the
circumferential order to an integer, $n\in\mathbb{N}_0$. For $n\ge1$ the
$\cos n\theta$ and $\sin n\theta$ solutions are degenerate and share a
frequency; for $n=0$ the mode is a singleton. Rigid-body motions on a
completely free ring are $n=0$, $s=0$ (translation) and $n=1$, $s=0$
(rotation). They are excluded from the elastic tables here, as in both
companion papers.

The radial operator is the $4\times4$ characteristic matrix already assembled
for the sector at a fixed separation constant. Its rows are, in order,
\begin{equation}
0=M_r(\text{inner}),\quad 1=M_r(\text{outer}),\quad
2=V_r(\text{inner}),\quad 3=V_r(\text{outer}),
\label{eq:rowmap}
\end{equation}
and searching $\det L(n,\Omega)=0$ at fixed integer $n$ \emph{is} the closed
ring: no $\theta$-edge condition is weakly enforced, so no residual is
generated and no residual bar is meaningful. The four Frobenius branches span
$\{J_n,Y_n,I_n,K_n\}$ under the radial map used throughout the series
(Supplementary Material, \S{}S.1).

Boundary conditions enter as row swaps on \eqref{eq:rowmap}:
\begin{itemize}\itemsep2pt
\item \textbf{Free (F).} Leave that edge's $M_r$ and $V_r$ rows in place.
\item \textbf{Clamped (C).} Replace \emph{both} of that edge's rows by $W$ and
      $W'$ evaluated at the same radius.
\item \textbf{Simply supported (S).} Replace only the $V_r$ row by $W$, leaving
      $M_r$.
\end{itemize}
Two conventions must be kept straight. Leissa's table titles list the
\emph{outer} edge first, while the interface used here is
$(\mathrm{bc}_{\text{inner}},\mathrm{bc}_{\text{outer}})$; SP-160
Table~2.22 (``Clamped, Free'') is therefore inner~F / outer~C, and
Tables~2.30 and~2.31 (``Free, Clamped'') are inner~C / outer~F. All nine
Vogel combinations are production tables below. Clamped--clamped was also
checked against the thin-ring strip limit
$\lamsq\to22.3733\,(a/w)^2$ of width $w$, giving ratios $0.9999$, $1.0000$
and $1.0000$ at $\ba=0.90$, $0.94$ and $0.97$, before the SP-160 comparison.

\subsection{Solid disk, out-of-plane}\label{sec:disk}
Boundedness at $r=0$ removes $Y_n$ and $K_n$. Two surviving coefficients and
two outer-edge conditions give a $2\times2$ characteristic matrix. The
practical difficulty is that the four series columns are mid-radius Taylor
initial conditions rather than the Bessel pair itself, and a truncated series
is finite at the origin, so ``smallest $|W|$ at $r=0$'' does not isolate the
regular branches. Regular solutions behave as $r^{|n|}$, which supplies two
homogeneous conditions --- an inner jet --- whose two-dimensional kernel
\emph{is} the regular pair:
\begin{equation}
n=0:\; W'=W'''=0;\qquad
n=1:\; W=W''=0;\qquad
n\ge2:\; W=W'=0.
\label{eq:jet}
\end{equation}
Projecting the two outer rows of the $4\times4$ onto that kernel gives the
production $2\times2$ (\S{}S.2).

A documented failure is worth naming, because it is the natural thing to try.
Evaluating the \emph{$4\times4$} at $R_i=0$ and looking for its zeros gives a
matrix of rank~2 at $n=0$ but rank~3 --- not~2 --- for $n\ge1$
(Table~\ref{tab:ranks}). Frequencies extracted from that matrix are not
Leissa's disk. The production path is the $2\times2$. A tiny-hole $4\times4$
at $\ba=10^{-3}$ is a check on that path, not the path itself, and it is a
check the implementation fails on two of four rows against a pre-registered
$10^{-3}$ relative bar (\S\ref{sec:disk-results}). That bar is not loosened.

\subsection{In-plane ring}\label{sec:ipring}
The in-plane characteristic matrix is $N_r$ and $N_{r\theta}$ at both arcs,
and the same substitution applies: fix $\xi=n$ at an integer, search $\Omega$.
There is no Kirchhoff corner term to reduce. The in-plane frequency parameter
of Irie, Yamada and Muramoto is defined from the extensional rigidity
$D_{\mathrm{ext}}=Eh/(1-\nu^2)$ as
$\lambda^2=\rho h a^2\omega^2/D_{\mathrm{ext}}$, that is
\begin{equation}
\lambda=\omega a\sqrt{\rho(1-\nu^{2})/E},
\label{eq:irielam}
\end{equation}
which is a different symbol from the flexural $\lamsq$ of
\eqref{eq:lam2} and must not reuse the flexural conversion factor. Their
``circular'' column is $\beta=b/a=0.01$, a tiny-hole control, and is not
claimed here as $R_i=0$.

\subsection{Conversion}\label{sec:conv}
All flexural comparisons use
\begin{equation}
\lamsq=\omega a^{2}\sqrt{\rho/D},\qquad a=R_o,\qquad
D=\frac{Eh^{3}}{12(1-\nu^{2})},
\label{eq:lam2}
\end{equation}
with $h$ the full thickness. The map from the solver's native $\Omega$ to
\eqref{eq:lam2} is recomputed from the geometry and material of each run and
is \emph{not} a universal constant. On the free--free steel plate used as the
fidelity anchor in \cite{mccune2026annular} ($R_o=8$~m, thickness $0.08$~m,
$\ba=0.5$) the factor happens to be $4\pi^{2}\approx39.478$; at $\ba=0.3$ it
is a different number, and reusing it is an error the series has made once and
now regression-tests against. The unit test is the anchor's mode~7: native
$\Omega=1.369611$ converts to $\lamsq=54.0755$, matching the finite-element
value to $0.01\%$. Native $\Omega$ is never printed with a $\lamsq$ label.
The full conversion algebra, including the separate in-plane map, is
Supplementary Material, \S{}S.3.

Poisson ratios are not mixed in a single comparison. Table~2.34 states
$\nu=1/3$ in its caption and is compared only at that value. Narita and the
finite-element models use the $\nu$ printed in their own captions ($0.30$ or
$0.33$); Irie's table is $\nu=0.3$. Vogel's comprehensive tables (2.18, 2.20, 2.22, 2.24, 2.26, 2.28, 2.30,
2.33, 2.35) print no Poisson ratio, so every sweep that touches those
tables was run at both $\nu=0.30$ and $\nu=1/3$. The paper-facing
comparison is the $\nu=0.30$ column; those tables' $\nu$ is not asserted
to be the $1/3$ of Table~2.34. Table~2.16 states $\nu=1/3$ and is
compared only at that value.

\section{Screen}\label{sec:screen}
The depth-and-gap residual bar used for completely free \emph{sectors}
\cite{mccune2026annular} does not travel to this geometry: applied unchanged
it rejects five genuine full-ring roots. The ring and disk therefore use a
different instrument, fixed before the production sweeps were launched:

\begin{enumerate}\itemsep2pt
\item A coarse scan of $\log_{10}|\det L|$ on a window local to each
      literature target, \emph{and} a gap-free coarse scan per $n$. Both, not
      either: the free--free wells at the anchor geometry are narrower than
      $5\times10^{-4}$ in $\Omega$, and a target-local window alone once
      missed a whole band.
\item Golden-section polish of each candidate minimum.
\item A real-part sign-flip ladder at
      $\delta=10^{-2},10^{-3},10^{-4},10^{-5},10^{-6}$ about the polished
      root. CLEAN means all five inner flips are present. FLOOR means the only
      inner miss is at $\delta=10^{-6}$ with all five outer flips present and
      the residual shrinking --- the arbitrary-precision working floor at
      $40$ digits, already named in \cite{mccune2026annular}. NOT means any
      inner miss at $\delta\ge10^{-5}$, or a broken outer/shrinking pattern.
\item Positive controls must come back CLEAN and negative controls must come
      back NOT, evaluated in the same run as the production rows.
\end{enumerate}

Scoring against a published number is separate from the ladder and is applied
to the relative miss alone: MATCH below $1\%$, WEAK below $3\%$, MISS
otherwise. FLOOR never converts a MATCH into a MISS. The ladder is never
retuned to include a published value, and the $1\%/3\%$ thresholds are not
adjusted table by table. This screen is neither the persistence screen of
\cite{mccune2026rect} nor the sector subdomain instrument of
\cite{mccune2026annular}; those names do not apply here.

\paragraph{A search window is a hypothesis.}
If $\log|\det L|$ falls monotonically to an endpoint of the scanned window,
the minimiser can return that endpoint. The returned value is then a property
of the window, not of the plate: it is $n$-independent and meaningless. The
sign-flip ladder catches this --- every such case in this work came back NOT
--- but the MATCH/WEAK/MISS gate above, computed from the relative miss alone,
does not. A post-hoc label (EXACT if the returned $\Omega$ is a window
endpoint, NEAR if within one coarse step of one, NONE otherwise) is therefore
attached to every result. The label only labels: it never moves the returned
$\Omega$, so no computed frequency changes and the solver's checkpoint version
key is unaffected. An audit of all $342$ stored rows of this study found $12$
endpoint returns and no near-misses, every one of them in the mixed-edge sweep
of \S\ref{sec:mixed}, which is the only sweep combining a steep
$\mathrm{d}\lamsq/\mathrm{d}(\ba)$ with a fixed absolute half-window. Those
cells are re-searched on a wide window where a root exists, and are marked in
the tables below. Endpoint values are not frequencies and are not printed as
such anywhere in this paper.

\section{Out-of-plane free--free annulus}\label{sec:ff-annulus}

\subsection{SP-160 Table 2.35}\label{sec:t235}
Table~2.35 of \cite{leissaSP160}, from Vogel and Skinner \cite{vogelSkinner},
gives eight $(n,s)$ rows at five radius ratios. Every one of the $80$ cells
across both Poisson ratios is MATCH or WEAK. At $\nu=0.30$
(Table~\ref{tab:t235-nu30}) all $40$ cells are MATCH, the largest miss being
$0.454\%$ at $(1,1)$, $\ba=0.1$, against a three-figure printed value. At
$\nu=1/3$ (Supplementary Material, \S{}S.4) $20$ cells are MATCH and $20$ are
WEAK, the largest miss $2.220\%$ at $(2,1)$, $\ba=0.9$. The WEAK cells at
$\nu=1/3$ cluster on the $(2,0)$ and $(3,0)$ branches, and those printed
values coincide with SP-160 Table~2.34, which is Raju's $\nu=1/3$ set
\cite{raju}, not Vogel's Table~2.35. Since Table~2.35 prints no Poisson ratio,
the paper-facing comparison for this table is the $\nu=0.30$ column, and
Table~2.35's $\nu$ is not asserted to be the $1/3$ of Table~2.34.

\begin{table}[htbp]\centering\footnotesize
\caption{Free--free annulus, $\nu=0.30$. Printed values are
SP-160 Table~2.35 \cite{leissaSP160,vogelSkinner}; solver values are
$\lamsq$ from $\det L(n,\Omega)=0$ on the $4\times4$ of
\S\ref{sec:ring}. All $40$ cells MATCH ($<1\%$); largest miss $0.454\%$.
Rigid-body $n=0,1$ at $s=0$ are excluded.}
\label{tab:t235-nu30}
\begin{tabular}{lcccccccccc}
\toprule
& \multicolumn{2}{c}{$\ba=0.1$} & \multicolumn{2}{c}{$\ba=0.3$}
& \multicolumn{2}{c}{$\ba=0.5$} & \multicolumn{2}{c}{$\ba=0.7$}
& \multicolumn{2}{c}{$\ba=0.9$}\\
\cmidrule(lr){2-3}\cmidrule(lr){4-5}\cmidrule(lr){6-7}\cmidrule(lr){8-9}\cmidrule(lr){10-11}
$(n,s)$ & printed & solver & printed & solver & printed & solver
& printed & solver & printed & solver\\
\midrule
$(2,0)$ & 5.30 & 5.3039 & 4.91 & 4.9060 & 4.28 & 4.2711 & 3.57 & 3.5725 & 2.94 & 2.9351 \\
$(3,0)$ & 12.4 & 12.437 & 12.26 & 12.266 & 11.4 & 11.425 & 9.86 & 9.8589 & 8.14 & 8.1449 \\
$(0,1)$ & 8.77 & 8.7747 & 8.36 & 8.3534 & 9.32 & 9.3135 & 13.2 & 13.163 & 34.9 & 34.834 \\
$(1,1)$ & 20.5 & 20.407 & 18.3 & 18.292 & 17.2 & 17.198 & 22.0 & 21.914 & 55.7 & 55.719 \\
$(2,1)$ & 34.9 & 34.930 & 33.0 & 32.973 & 31.1 & 31.115 & 37.8 & 37.841 & 93.8 & 93.794 \\
$(3,1)$ & 53.0 & 52.965 & 51.0 & 51.005 & 47.4 & 47.462 & 55.7 & 55.701 & 135.0 & 135.41 \\
$(0,2)$ & 38.2 & 38.236 & 50.4 & 50.353 & 92.3 & 92.308 & 251.0 & 250.58 & 2238.0 & 2238.9 \\
$(1,2)$ & 59.0 & 59.072 & 58.8 & 58.784 & 96.3 & 96.266 & 253.0 & 253.13 & 2240.0 & 2240.8 \\
\bottomrule
\end{tabular}
\end{table}

\subsection{Narita's five-figure hold-out}\label{sec:narita2a}
Table~2.35 is a three-figure table, so it cannot discriminate at the level the
determinant actually resolves. Narita's Table~2(a) \cite{narita1984} at
$\ba=0.3$, $\nu=0.3$ is the hold-out: four five-figure values that were not
used in any development decision. All four are recovered to between
$3.7\times10^{-6}$ and $1.6\times10^{-5}$ relative
(Table~\ref{tab:narita2a}); the printed five figures hold in every case. The
$(3,0)$ row is FLOOR rather than CLEAN --- its only inner miss is at
$\delta=10^{-6}$, with all five outer flips present and shrinking --- which is
the precision floor of the ladder, not a missing frequency.

\begin{table}[htbp]\centering\small
\caption{Narita Table~2(a) \cite{narita1984} five-figure hold-out,
free--free annulus, $b/a=0.3$, $\nu=0.3$.}
\label{tab:narita2a}
\begin{tabular}{lcccc}
\toprule
$(n,s)$ & printed $\lamsq$ & solver $\lamsq$ & rel.\ miss & ladder\\
\midrule
$(2,0)$ & 4.9060 & 4.9060 & 6.4e-06 & CLEAN \\
$(0,1)$ & 8.3535 & 8.3535 & 3.7e-06 & CLEAN \\
$(3,0)$ & 12.266 & 12.266 & 1.6e-05 & FLOOR \\
$(1,1)$ & 18.292 & 18.292 & 1.1e-05 & CLEAN \\
\bottomrule
\end{tabular}
\end{table}

\subsection{Polar-orthotropic free--free hold-out}\label{sec:narita3}
The same $4\times4$ accepts a polar-orthotropic material with no new
row map. Narita's three parameters $[D_\theta/D_r,\,H/D_r,\,\nu_\theta]$
map onto the radial/hoop moduli, the torsional combination $T$, and the
reciprocal Poisson pair already used on the sector; the isotropic
reduction ($E_r=E_\theta$) reproduces the isotropic $\log|\det L|$ to
machine precision at $n=0$ and $n=2$. Table~3 of \cite{narita1984} is
then a data-entry check. The high-modulus graphite-epoxy plate
($[0.04,\,0.055,\,0.012]$) at $\ba=0.5$ and $\ba=0.3$ is recovered at
$10/10$ MATCH (Table~\ref{tab:narita3}), the largest miss
$4.1\times10^{-4}$. The one NOT ladder, $(1,1)$ at $\ba=0.3$, is
frequency-true at $6.9\times10^{-6}$ and is not retuned.

\begin{table}[htbp]\centering\small
\caption{Narita Table~3 \cite{narita1984}, high-modulus graphite epoxy
$[D_\theta/D_r,\,H/D_r,\,\nu_\theta]=[0.04,\,0.055,\,0.012]$, free--free
annulus. Narita's $\Omega=\omega a^{2}\sqrt{\rho h/D_r}$ is this paper's
$\lamsq$ when $D=D_r$ and $h$ is the full thickness. All ten cells MATCH.}
\label{tab:narita3}
\begin{tabular}{lccccc}
\toprule
$\ba$ & $(n,s)$ & printed $\Omega$ & solver $\lamsq$ & rel.\ miss & ladder\\
\midrule
0.5 & $(2,0)$ & 0.935 & 0.9348 & 2.22e-04 & FLOOR \\
0.5 & $(0,1)$ & 1.918 & 1.9174 & 3.03e-04 & FLOOR \\
0.5 & $(3,0)$ & 2.503 & 2.5029 & 2.4e-05 & CLEAN \\
0.5 & $(4,0)$ & 4.636 & 4.6363 & 5.6e-05 & FLOOR \\
0.5 & $(1,1)$ & 3.961 & 3.9614 & 9.7e-05 & FLOOR \\
0.3 & $(2,0)$ & 1.132 & 1.1315 & 4.06e-04 & CLEAN \\
0.3 & $(0,1)$ & 1.693 & 1.6930 & 1.8e-05 & CLEAN \\
0.3 & $(3,0)$ & 2.897 & 2.8971 & 4.6e-05 & FLOOR \\
0.3 & $(4,0)$ & 5.180 & 5.1795 & 9.6e-05 & FLOOR \\
0.3 & $(1,1)$ & 4.519 & 4.5190 & 6.9e-06 & NOT \\
\bottomrule
\end{tabular}
\end{table}

\subsection{Fidelity anchors and finite elements}\label{sec:ff-fe}
Five free--free roots at $\ba=0.5$, $\nu=0.30$ are carried forward from
\cite{mccune2026annular} as a fidelity gate on the production interface rather
than as a new result: the ring path must reproduce them in native $\Omega$
before any table is believed. It does, to between $2.9\times10^{-6}$ and
$1.4\times10^{-5}$ (Table~\ref{tab:five}).

Independent cyclic-symmetry SHELL281 models at the same geometry give the
elastic listing $4.2645$, $9.3114$, $11.4082$, $17.0973$, $21.0296$, agreeing
with the solver to between $0.023\%$ and $0.591\%$. Mesh convergence at that
geometry is $4.264506$, $4.264461$, $4.264453$ for $64$, $96$ and $128$
circumferential divisions --- a relative change of $10^{-5}$ between the two
finest meshes. Contour plots of $U_Z$ confirm the published $(n,s)$
assignment on every mode in the plotted set, including the $(3,0)$, $(1,1)$
and $(4,0)$ rows whose labels had previously been inferred from frequency
ordering alone. At $\ba=0.9$ the finite-element modes $15$--$16$ are in-plane
($|U_Z|\sim10^{-13}$) and must not be given a flexural $n$; the flexural
$(0,1)$ mode at that geometry is the first axisymmetric elastic mode above
that pair.

\begin{table}[htbp]\centering\small
\caption{Free--free annulus at $\ba=0.5$, $\nu=0.30$: the five fidelity
anchors of \cite{mccune2026annular} in native $\Omega$, and the same modes
against independent SHELL281 \cite{ansys} in $\lamsq$. Ladder FLOOR here is
the $40$-digit working floor, not a missing root.}
\label{tab:five}
\begin{tabular}{lccccccc}
\toprule
$(n,s)$ & anchor $\Omega$ & solver $\Omega$ & rel. & ladder
& FE $\lamsq$ & solver $\lamsq$ & rel.\ (\%)\\
\midrule
$(2,0)$ & 0.1081885366 & 0.1081900894 & 1.4e-05 & FLOOR & 4.2645 & 4.2711 & 0.155 \\
$(0,1)$ & 0.2359132157 & 0.2359147685 & 6.6e-06 & FLOOR & 9.3114 & 9.3135 & 0.023 \\
$(3,0)$ & 0.2894098129 & 0.2894082601 & 5.4e-06 & FLOOR & 11.4082 & 11.4254 & 0.151 \\
$(1,1)$ & 0.4356362253 & 0.4356346725 & 3.6e-06 & FLOOR & 17.0973 & 17.1983 & 0.591 \\
$(4,0)$ & 0.5336384799 & 0.5336400327 & 2.9e-06 & FLOOR & 21.0296 & 21.067 & 0.178 \\
\bottomrule
\end{tabular}
\end{table}

\FloatBarrier
\section{Out-of-plane free disk}\label{sec:disk-results}
The solid disk uses the $2\times2$ of \S\ref{sec:disk}. Narita's Table~2(b)
\cite{narita1984} at $\nu=0.33$ is again the five-figure hold-out and is
recovered $4/4$ at between $1.7\times10^{-6}$ and $9.4\times10^{-6}$ relative
(Table~\ref{tab:narita2b}). The $(2,0)$ hold-out is MATCH on the relative miss
and NOT on the sign-flip ladder; it is trusted because the finite elements sit
on the same value ($5.2584$ against the solver's $5.2620$, $0.069\%$). SP-160
Table~2.5 \cite{leissaSP160}, six rows at the same Poisson ratio, is recovered
$6/6$ (Table~\ref{tab:leissa25}); its largest disagreement is the
$(2,0)$ row, printed $5.253$ against the solver's $5.262$, a $0.172\%$
difference that sits inside the $1\%$ bar appropriate to a four-figure table
and is resolved in the solver's favour by the five-figure hold-out and by the
finite elements, both of which give $5.262$ and $5.2584$ respectively. The
Table~2.5 $(0,1)$ row is likewise MATCH on error and NOT on the ladder, and is
read the same way.

Independent SHELL281 models of the free disk at $\nu=0.33$ give $5.2584$,
$9.0679$, $12.2284$ and $20.5010$ for the four Narita modes: $0.068\%$,
$0.011\%$, $0.13\%$ and $0.058\%$ from the printed values and $0.069\%$,
$0.011\%$, $0.13\%$ and $0.059\%$ from the $2\times2$. Contours confirm the
published $(n,s)$ on all six Table~2.5 rows.

\begin{table}[htbp]\centering\small
\caption{Free solid disk, $\nu=0.33$, from the $2\times2$ of \S\ref{sec:disk},
against Narita Table~2(b) \cite{narita1984} (five-figure hold-out) and
independent SHELL281 \cite{ansys}.}
\label{tab:narita2b}
\begin{tabular}{lcccccc}
\toprule
$(n,s)$ & printed $\lamsq$ & solver $\lamsq$ & rel.\ miss & FE $\lamsq$
& solver vs FE (\%) & ladder\\
\midrule
$(2,0)$ & 5.262 & 5.2620 & 2.9e-06 & 5.2584 & 0.069 & NOT \\
$(0,1)$ & 9.0689 & 9.0689 & 1.7e-06 & 9.0679 & 0.011 & FLOOR \\
$(3,0)$ & 12.244 & 12.244 & 9.4e-06 & 12.2284 & 0.127 & CLEAN \\
$(1,1)$ & 20.513 & 20.513 & 4.7e-06 & 20.5010 & 0.059 & CLEAN \\
\bottomrule
\end{tabular}
\end{table}

\begin{table}[htbp]\centering\small
\caption{Free solid disk, $\nu=0.33$, against SP-160 Table~2.5
\cite{leissaSP160}. All six MATCH. The $(2,0)$ row is the largest
disagreement at $0.172\%$; the five-figure hold-out and the finite elements
both sit on the solver's value.}
\label{tab:leissa25}
\begin{tabular}{lccccc}
\toprule
$(n,s)$ & printed $\lamsq$ & solver $\lamsq$ & rel.\ (\%) & FE $\lamsq$
& ladder\\
\midrule
$(2,0)$ & 5.253 & 5.2620 & 0.172 & 5.2584 & CLEAN \\
$(3,0)$ & 12.23 & 12.244 & 0.114 & 12.2284 & CLEAN \\
$(4,0)$ & 21.6 & 21.527 & 0.337 & 21.4890 & CLEAN \\
$(0,1)$ & 9.084 & 9.0689 & 0.166 & 9.0679 & NOT \\
$(1,1)$ & 20.52 & 20.513 & 0.034 & 20.5010 & CLEAN \\
$(2,1)$ & 35.25 & 35.243 & 0.021 & 35.1988 & CLEAN \\
\bottomrule
\end{tabular}
\end{table}

\subsection{Why the \texorpdfstring{$4\times4$}{4x4} at
\texorpdfstring{$R_i=0$}{Ri=0} is not the disk}\label{sec:rank}
Substituting $R_i=0$ into the ring $4\times4$ and looking for its zeros is the
obvious shortcut and it is wrong. Table~\ref{tab:ranks} lists the ordered
$\log_{10}$ singular values of that matrix at a representative frequency: at
$n=0$ there are two small values and the rank is~2, but at $n=1,2,3$ only one
singular value collapses and the rank is~3. A rank-3 $4\times4$ has a
one-dimensional kernel that is not the two-dimensional regular subspace, and
its determinant zeros are not the plate's frequencies. This is the reason the
production disk path is the $2\times2$ and why that fact is enforced by a
regression test rather than by convention. The full regularity argument
behind this rank count is Supplementary Material, \S{}S.2.

\begin{table}[htbp]\centering\small
\caption{Ordered $\log_{10}$ singular values of the ring $4\times4$ evaluated
at $R_i=0$, with the observed and expected numerical rank. Rank~3 for
$n\ge1$ is why this matrix is not the disk.}
\label{tab:ranks}
\begin{tabular}{lcccccc}
\toprule
& $\sigma_1$ & $\sigma_2$ & $\sigma_3$ & $\sigma_4$ & rank & expected\\
\midrule
$n=0$ & 0.43 & -0.96 & -16.09 & -31.88 & 2 & 2 \\
$n=1$ & 0.60 & -2.41 & -2.88 & -15.51 & 3 & 3 \\
$n=2$ & 0.60 & -1.38 & -3.81 & -16.94 & 3 & 3 \\
$n=3$ & 0.60 & -1.89 & -4.91 & -16.61 & 3 & 3 \\
\bottomrule
\end{tabular}
\end{table}

\subsection{A named failure: the tiny-hole control}\label{sec:tinyhole}
A separate control asks whether the $4\times4$ at a very small hole,
$\ba=10^{-3}$, reproduces the $2\times2$ disk to a pre-registered $10^{-3}$
relative tolerance. It does not (Table~\ref{tab:tinyhole}): $(2,\cdot)$ and
$(0,\cdot)$ differ by $0.940\%$ and $0.545\%$, and the control is reported as
a FAIL. The hole roots are real and deep, and both failing rows have ladder
NOT while the two passing rows are CLEAN, so this is the residual physical
difference between a punctured and a solid plate at that hole size, not a
missed disk frequency. The tolerance is not loosened to make the control pass.

\begin{table}[htbp]\centering\small
\caption{Tiny-hole control, $\ba=10^{-3}$, $\nu=0.33$: the $4\times4$ at a
small hole against the production $2\times2$. Pre-registered bar $10^{-3}$
relative; the control is a named FAIL and the bar is not loosened.}
\label{tab:tinyhole}
\begin{tabular}{lccccc}
\toprule
mode & disk $2\times2$ & hole $4\times4$ & rel.\ (\%) & ladder & control\\
\midrule
$(2,\cdot)$ & 5.2620 & 5.3115 & 0.940 & NOT & \textbf{FAIL} \\
$(0,\cdot)$ & 9.0689 & 9.1183 & 0.545 & NOT & \textbf{FAIL} \\
$(3,\cdot)$ & 12.244 & 12.245 & 0.007 & CLEAN & PASS \\
$(1,\cdot)$ & 20.513 & 20.524 & 0.054 & CLEAN & PASS \\
\bottomrule
\end{tabular}
\end{table}

\FloatBarrier
\section{Mixed inner and outer edges}\label{sec:mixed}
Mixed inner/outer conditions are the content of this paper that is not a
free--free echo of the companion papers, and they are reached by the row swap
of \S\ref{sec:ring} with no new derivation. Because there is no $\theta$-edge,
all nine of Vogel and Skinner's inner/outer pairs are different row pairs of
one $4\times4$. Clamping one arc of a ring is also the mechanical boundary
condition of a stator disk; that observation is a statement about the
boundary condition only, and no electromechanical claim is made here.

Table~\ref{tab:nine} is the production scoreboard. Combined across both
Poisson ratios the nine tables are $704/710$ MATCH${}+{}$WEAK. Named source
skips (Table~2.30's $n=0$ column at $\ba=0.9$; Table~2.26's printed $933$;
Table~2.18's two blank cells at $\ba=0.9$) are excluded from the
denominators on source-internal grounds, not because the solver disagrees.
The $1\%/3\%$ bar is unchanged.

\begin{table}[htbp]\centering\small
\caption{Nine Vogel inner/outer combinations as row pairs of the
$4\times4$. Leissa titles list the outer edge first; the local interface is
(inner, outer). Scores are MATCH${}+{}$WEAK over both $\nu=0.30$ and
$\nu=1/3$. Blanks and named source skips are not in the denominator.}
\label{tab:nine}
\begin{tabular}{llccc}
\toprule
SP-160 & inner--outer & score & blanks / skips & open misses\\
\midrule
2.35 & F--F & $80/80$ & --- & --- \\
2.18 & C--C & $76/76$ & 4 blanks ($n=2$, $\ba=0.9$) & --- \\
2.20 & S--C & $80/80$ & --- & --- \\
2.22 & F--C & $80/80$ & --- & --- \\
2.24 & C--S & $80/80$ & --- & --- \\
2.26 & S--S & $78/78$ & printed $933$ at $(2,1)$, $\ba=0.3$ & --- \\
2.28 & F--S & $80/80$ & --- & --- \\
2.30 & C--F & $72/76$ & $n=0$ at $\ba=0.9$ & $(1,0)$ at $\ba=0.1$, $0.3$ \\
2.33 & S--F & $78/80$ & --- & $(1,0)$ at $\ba=0.1$ \\
\midrule
combined & & $704/710$ & & \\
\bottomrule
\end{tabular}
\end{table}

\subsection{Vogel and Skinner, Tables 2.22 and 2.30}\label{sec:t222}
Table~2.22 (inner~F, outer~C) is recovered at $80/80$ across both Poisson
ratios and all five radius ratios including $\ba=0.9$, with $78$ cells MATCH
and $2$ WEAK (Table~\ref{tab:t222-nu30}; the $\nu=1/3$ half is \S{}S.4). The
thin-ring column is worth naming explicitly: the $4\times4$ does not degrade
at $\ba=0.9$, where the printed values run to $6183$ and the solver returns
$6185.1$.

\begin{table}[htbp]\centering\footnotesize
\caption{Mixed edges, inner free / outer clamped, $\nu=0.30$. Printed values
are SP-160 Table~2.22 \cite{leissaSP160,vogelSkinner} (Leissa's title lists
the outer edge first). 80/80 MATCH${}+{}$WEAK over both Poisson ratios.}
\label{tab:t222-nu30}
\begin{tabular}{lcccccccccc}
\toprule
& \multicolumn{2}{c}{$\ba=0.1$} & \multicolumn{2}{c}{$\ba=0.3$}
& \multicolumn{2}{c}{$\ba=0.5$} & \multicolumn{2}{c}{$\ba=0.7$}
& \multicolumn{2}{c}{$\ba=0.9$}\\
\cmidrule(lr){2-3}\cmidrule(lr){4-5}\cmidrule(lr){6-7}\cmidrule(lr){8-9}\cmidrule(lr){10-11}
$(n,s)$ & printed & solver & printed & solver & printed & solver
& printed & solver & printed & solver\\
\midrule
$(0,0)$ & 10.2 & 10.159 & 11.4 & 11.424 & 17.7 & 17.715 & 43.1 & 43.142 & 360 & 360.35 \\
$(1,0)$ & 21.1 & 21.195 & 19.5 & 19.540 & 22.0 & 22.015 & 45.3 & 45.333 & 362 & 361.50 \\
$(2,0)$ & 34.5 & 34.535 & 32.5 & 32.594 & 32.0 & 32.115 & 51.5 & 51.585 & 365 & 364.96 \\
$(3,0)$ & 51.0 & 50.989 & 49.1 & 49.069 & 45.8 & 45.812 & 61.3 & 61.289 & 370 & 370.70 \\
$(0,1)$ & 39.5 & 39.521 & 51.7 & 51.745 & 93.8 & 93.847 & 253 & 251.64 & 2219 & 2219.5 \\
$(1,1)$ & 60.0 & 60.061 & 59.8 & 59.759 & 97.3 & 97.376 & 254 & 253.72 & 2220 & 2220.9 \\
$(2,1)$ & 83.4 & 83.478 & 79.0 & 79.061 & 108.0 & 107.49 & 259 & 259.91 & 2225 & 2225.2 \\
$(0,2)$ & 90.4 & 90.445 & 132.0 & 132.41 & 253.0 & 252.19 & 692 & 692.01 & 6183 & 6185.1 \\
\bottomrule
\end{tabular}
\end{table}

Table~2.30 (inner~C, outer~F) is reported here as $72/76$
(Table~\ref{tab:t230-nu30}). Four cells are excluded as one named source
defect and four are direct misses. Both need justifying.

The excluded cells are the whole $n=0$ column at $\ba=0.9$, at both Poisson
ratios. At that radius ratio the thin-ring $n=0$ and $n=1$ branches are nearly
degenerate, and Vogel's \emph{own} Table~2.22 at the same $\ba$ prints that
degeneracy: $360$ against $362$ for $s=0$, and $2219$ against $2220$ for
$s=1$. Table~2.30 at the same $\ba$ instead prints $51.5$ against $345$ and
$970$ against $2189$. The $n=1$ entries agree with the $4\times4$ to $0.06\%$;
the $n=0$ entries agree with nothing, and no confident root exists anywhere in
$\lamsq\in[15,155]$ at the first of them. The argument for excluding that
column is therefore internal to the source --- Vogel's other table plus the
physical degeneracy --- and not ``the solver disagrees''. The MATCH/WEAK bar
is unchanged. For the record, the production sweep scored this table $72/78$,
excluding only the printed $51.5$; the paper's denominator differs from the
sweep's by treating the column as one skip rather than one cell.

The four direct misses are the $(1,0)$ entries at $\ba=0.1$ and $\ba=0.3$, at
both Poisson ratios: printed $3.14$ against solver $3.478$ ($10.8\%$) and
printed $6.33$ against solver $6.552$ ($3.5\%$). They are named, not dropped, and
they remain open: both lie on the $n=1$, $s=0$ branch at the two smallest
radius ratios, where the printed values are also the smallest and coarsest in
the table.
Those two tables alone are 152/156; the nine-combination total is 704/710.

\begin{table}[htbp]\centering\footnotesize
\caption{Mixed edges, inner clamped / outer free, $\nu=0.30$. Printed values
are SP-160 Table~2.30 \cite{leissaSP160,vogelSkinner}. $^{\dagger}$~the $n=0$
column at $b/a=0.9$ is excluded as one named source defect, on the
source-internal grounds given in the text. \winend~the production search
returned its window endpoint at this cell; an endpoint is not a frequency
and is not printed (see \S\ref{sec:screen}). The re-searched $(1,1)$ root at
$b/a=0.1$ is WEAK at $1.37\%$ ($\nu=0.30$) and $1.56\%$ ($\nu=1/3$) and is
scored WEAK, as the stored value was. Paper score 72/76.}
\label{tab:t230-nu30}
\begin{tabular}{lcccccccccc}
\toprule
& \multicolumn{2}{c}{$\ba=0.1$} & \multicolumn{2}{c}{$\ba=0.3$}
& \multicolumn{2}{c}{$\ba=0.5$} & \multicolumn{2}{c}{$\ba=0.7$}
& \multicolumn{2}{c}{$\ba=0.9$}\\
\cmidrule(lr){2-3}\cmidrule(lr){4-5}\cmidrule(lr){6-7}\cmidrule(lr){8-9}\cmidrule(lr){10-11}
$(n,s)$ & printed & solver & printed & solver & printed & solver
& printed & solver & printed & solver\\
\midrule
$(1,0)$ & 3.14 & 3.4776 & 6.33 & 6.5523 & 13.3 & 13.290 & 37.5 & 37.498 & 345 & 345.17 \\
$(0,0)$ & 4.23 & 4.2373 & 6.66 & 6.6604 & 13.0 & 13.024 & 37.0 & 36.953 & 51.5$^{\dagger}$ & \winend \\
$(2,0)$ & 5.62 & 5.6202 & 7.96 & 7.9565 & 14.7 & 14.704 & 39.3 & 39.277 & 347 & 347.47 \\
$(3,0)$ & 12.4 & 12.448 & 13.27 & 13.276 & 18.5 & 18.562 & 42.6 & 42.655 & 352 & 351.32 \\
$(0,1)$ & 25.3 & 25.262 & 42.6 & 42.614 & 85.1 & 85.033 & 239 & 239.84 & 970$^{\dagger}$ & \winend \\
$(1,1)$ & 27.3 & \winend & 44.6 & 44.631 & 86.7 & 86.706 & 241 & 241.28 & 2189 & 2190.1 \\
$(2,1)$ & 37.0 & 36.943 & 50.9 & 50.947 & 91.7 & 91.738 & 246 & 245.59 & 2194 & 2193.9 \\
$(3,1)$ & 53.2 & 53.207 & 62.1 & 62.053 & 100 & 100.17 & 253 & 252.73 & 2200 & 2200.3 \\
\bottomrule
\end{tabular}
\end{table}

\subsection{The axisymmetric exact values}\label{sec:t216}
SP-160 Table~2.16 collects exact axisymmetric frequencies at $\ba=0.5$,
$\nu=1/3$, for every inner/outer pair except free--free. All eight printed
cells are MATCH (Table~\ref{tab:t216}), including clamped--clamped at
printed $89.30$ against solver $89.251$ ($0.055\%$). These are three- and
four-figure printed values and the agreement is well inside the bar.

\begin{table}[htbp]\centering\small
\caption{Exact axisymmetric values, SP-160 Table~2.16
\cite{leissaSP160}, $\ba=0.5$, $\nu=1/3$, $n=0$. All eight printed
combinations MATCH. Free--free is omitted in the source.}
\label{tab:t216}
\begin{tabular}{lcccc}
\toprule
inner--outer & printed $\lamsq$ & solver $\lamsq$ & rel.\ (\%) & ladder\\
\midrule
C--C & 89.30 & 89.2507 & 0.055 & CLEAN \\
S--C & 64.06 & 63.8563 & 0.318 & CLEAN \\
F--C & 17.51 & 17.5980 & 0.503 & CLEAN \\
C--S & 59.91 & 59.8987 & 0.019 & FLOOR \\
S--S & 40.01 & 40.0111 & 0.003 & FLOOR \\
F--S & 5.040 & 5.0432 & 0.063 & FLOOR \\
C--F & 13.05 & 13.0894 & 0.302 & CLEAN \\
S--F & 4.060 & 4.0700 & 0.247 & CLEAN \\
\bottomrule
\end{tabular}
\end{table}

\subsection{Clamped--clamped and simply supported--simply supported}
\label{sec:t218}
Table~2.18 (inner~C, outer~C) is $76/76$ across both Poisson ratios
(Table~\ref{tab:t218-nu30}; $\nu=1/3$ is identical in $\lamsq$ to the
precision printed, as expected of a fully clamped Kirchhoff plate). The
source leaves $(2,0)$ and $(2,1)$ at $\ba=0.9$ blank; those four cells
(two $\nu$) are not invented. The thin-ring column again holds: printed
$2237$ against $2237.2$.

\begin{table}[htbp]\centering\footnotesize
\caption{Clamped--clamped annulus, $\nu=0.30$. Printed values are
SP-160 Table~2.18 \cite{leissaSP160,vogelSkinner}. 76/76
MATCH${}+{}$WEAK over both Poisson ratios. An em-dash marks a blank in
the source; not a miss.}
\label{tab:t218-nu30}
\begin{tabular}{lcccccccccc}
\toprule
& \multicolumn{2}{c}{$\ba=0.1$} & \multicolumn{2}{c}{$\ba=0.3$}
& \multicolumn{2}{c}{$\ba=0.5$} & \multicolumn{2}{c}{$\ba=0.7$}
& \multicolumn{2}{c}{$\ba=0.9$}\\
\cmidrule(lr){2-3}\cmidrule(lr){4-5}\cmidrule(lr){6-7}\cmidrule(lr){8-9}\cmidrule(lr){10-11}
$(n,s)$ & printed & solver & printed & solver & printed & solver
& printed & solver & printed & solver\\
\midrule
$(0,0)$ & 27.3 & 27.281 & 45.2 & 45.346 & 89.2 & 89.251 & 248 & 248.40 & 2237 & 2237.2 \\
$(1,0)$ & 28.4 & 28.766 & 46.6 & 46.644 & 90.2 & 90.230 & 249 & 249.16 & 2238 & 2237.8 \\
$(2,0)$ & 36.7 & 36.617 & 51.0 & 51.139 & 93.3 & 93.321 & 251 & 251.48 & --- & --- \\
$(3,0)$ & 51.2 & 51.219 & 60.0 & 60.033 & 99.0 & 98.928 & 256 & 255.44 & 2243 & 2242.7 \\
$(0,1)$ & 75.3 & 75.366 & 125 & 125.36 & 246 & 246.34 & 686 & 684.99 & 6167 & 6167.1 \\
$(1,1)$ & 78.6 & 78.635 & 127 & 127.38 & 248 & 247.74 & 686 & 686.04 & 6167 & 6167.9 \\
$(2,1)$ & 90.5 & 90.448 & 134 & 133.67 & 253 & 251.97 & 689 & 689.19 & --- & --- \\
$(3,1)$ & 112 & 112.11 & 145 & 144.83 & 259 & 259.16 & 694 & 694.47 & 6174 & 6174.5 \\
\bottomrule
\end{tabular}
\end{table}

Table~2.26 (inner~S, outer~S) is $78/78$ (Table~\ref{tab:t226-nu30}).
The source prints $933$ at $(2,1)$, $\ba=0.3$. Neighbours on that row are
$71.7$ and $167$; a search centred on $93.3$ returns $\lamsq=93.371$
(CLEAN, one interior radial node). A search centred on the printed $933$
returns a sixth-radial-node root at $933.60$. The printed cell is a lost
decimal, excluded as one named typeset at both Poisson ratios, and the
MATCH/WEAK bar is not retuned to absorb $933$.

\begin{table}[htbp]\centering\footnotesize
\caption{Simply supported--simply supported annulus, $\nu=0.30$. Printed
values are SP-160 Table~2.26 \cite{leissaSP160,vogelSkinner}. 78/78
MATCH${}+{}$WEAK over both Poisson ratios.
$^{\dagger}$~printed $933$ is a lost decimal; the solver value is the
$s=1$ root from a search at $93.3$.}
\label{tab:t226-nu30}
\begin{tabular}{lcccccccccc}
\toprule
& \multicolumn{2}{c}{$\ba=0.1$} & \multicolumn{2}{c}{$\ba=0.3$}
& \multicolumn{2}{c}{$\ba=0.5$} & \multicolumn{2}{c}{$\ba=0.7$}
& \multicolumn{2}{c}{$\ba=0.9$}\\
\cmidrule(lr){2-3}\cmidrule(lr){4-5}\cmidrule(lr){6-7}\cmidrule(lr){8-9}\cmidrule(lr){10-11}
$(n,s)$ & printed & solver & printed & solver & printed & solver
& printed & solver & printed & solver\\
\midrule
$(0,0)$ & 14.5 & 14.485 & 21.1 & 21.079 & 40.0 & 40.043 & 110 & 110.06 & 988 & 987.27 \\
$(1,0)$ & 16.7 & 16.776 & 23.3 & 23.317 & 41.8 & 41.797 & 112 & 111.44 & 988 & 988.38 \\
$(2,0)$ & 25.9 & 25.936 & 30.2 & 30.273 & 47.1 & 47.089 & 116 & 115.59 & 993 & 991.70 \\
$(3,0)$ & 40.0 & 39.976 & 42.0 & 41.910 & 56.0 & 55.957 & 122 & 122.49 & 998 & 997.24 \\
$(0,1)$ & 51.7 & 51.781 & 81.8 & 81.737 & 159 & 158.64 & 439 & 439.18 & 3948 & 3948.3 \\
$(1,1)$ & 56.5 & 56.507 & 84.6 & 84.635 & 161 & 160.57 & 441 & 440.59 & 3948 & 3949.4 \\
$(2,1)$ & 71.7 & 71.686 & 933$^{\dagger}$ & 93.371 & 167 & 166.35 & 444 & 444.84 & 3952 & 3952.7 \\
$(3,1)$ & 94.7 & 94.712 & 108 & 108.20 & 177 & 176.01 & 453 & 451.92 & 3958 & 3958.2 \\
\bottomrule
\end{tabular}
\end{table}

The four combinations that mix a simply-supported edge with free or
clamped (Tables~2.20, 2.24, 2.28, 2.33) are in Supplementary Material
\S{}S.4. Three of them are $80/80$. Table~2.33 (inner~S, outer~F) is
$78/80$: the two misses are one printed cell at both Poisson ratios,
$(1,0)$ at $\ba=0.1$, printed $2.30$ against solver $2.409$ ($\nu=1/3$,
$4.72\%$, CLEAN) and $2.438$ ($\nu=0.30$, $6.02\%$, FLOOR). That is the
real nearest root, not a window endpoint, and it is the same small-$\ba$,
$n=1$, $s=0$ class as the open misses of Table~2.30. The bar is not
loosened; the combination is not sent to finite elements.

\FloatBarrier
\subsection{Southwell's Table 2.31}\label{sec:southwell}
Southwell's values for the free-outer, clamped-inner annulus
\cite{southwell}, compiled as SP-160 Table~2.31 at $\nu=0.3$, are the hardest
comparison in this paper and are reported as such:
\textbf{$13/18$} under the same $1\%/3\%$ bar as every other table
(Table~\ref{tab:southwell}). No second, table-specific metric is substituted
for that score.

\begin{table}[htbp]\centering\small
\caption{Southwell's values \cite{southwell} as compiled in SP-160
Table~2.31 \cite{leissaSP160}, inner clamped / outer free, $\nu=0.3$, lowest
root per $n$. Score $13/18$ under the same bar as every other table.
$^{\ddagger}$~the production search returned its window endpoint at this cell;
the value shown is from a wide re-search and is CLEAN with $s=0$. All three
re-searched rows were, and remain, MISS, so the score is unchanged.}
\label{tab:southwell}
\begin{tabular}{cccccc}
\toprule
$n$ & $\ba$ & printed $\lamsq$ & solver $\lamsq$ & rel.\ (\%) & gate\\
\midrule
0 & 0.276 & 6.25 & 6.2426 & 0.119 & MATCH \\
0 & 0.642 & 25.0 & 25.745 & 2.981 & WEAK \\
0 & 0.84 & 81.0 & 132.99$^{\ddagger}$ & 64.185 & MISS \\
1 & 0.06 & 2.82 & 2.8148 & 0.184 & MATCH \\
1 & 0.397 & 9.00 & 9.0207 & 0.230 & MATCH \\
1 & 0.603 & 21.2 & 21.258 & 0.275 & MATCH \\
1 & 0.634 & 25.0 & 25.068 & 0.274 & MATCH \\
1 & 0.771 & 64.0 & 64.762 & 1.191 & WEAK \\
1 & 0.827 & 121.0 & 114.17 & 5.642 & MISS \\
2 & 0.186 & 6.25 & 6.2947 & 0.715 & MATCH \\
2 & 0.349 & 9.00 & 9.0294 & 0.327 & MATCH \\
2 & 0.522 & 16.0 & 16.011 & 0.072 & MATCH \\
2 & 0.769 & 64.0 & 65.583 & 2.473 & WEAK \\
2 & 0.81 & 100 & 96.526 & 3.474 & MISS \\
3 & 0.43 & 16.0 & 15.785 & 1.342 & WEAK \\
3 & 0.59 & 25.0 & 24.964 & 0.146 & MATCH \\
3 & 0.71 & 49.0 & 45.346$^{\ddagger}$ & 7.457 & MISS \\
3 & 0.82 & 100 & 111.04$^{\ddagger}$ & 11.041 & MISS \\
\bottomrule
\end{tabular}
\end{table}

\paragraph{The split is a sensitivity, not an accuracy.}
Southwell's method assumed round values of the Bessel argument and inverted
for the radius ratio, so his tabulated $\ba$ is a computed quantity printed to
two or three decimals while his $\lamsq$ is exact by construction. Holding the
printed $\lamsq$ fixed and root-finding on $\ba$ instead recovers a CLEAN,
$s=0$ root on \emph{all eighteen} rows. Every row then carries a radius-ratio
discrepancy of order $10^{-3}$ --- including the rows that match. For the
thirteen MATCH and WEAK rows that discrepancy spans
$1.6\times10^{-4}$ to $6.9\times10^{-3}$, a band that overlaps the misses.
What separates the two groups is not the size of the discrepancy but the local
sensitivity: the predicted forward error
\begin{equation}
\left(\frac{\mathrm{d}\lamsq}{\mathrm{d}(\ba)}\right)
\frac{|\delta(\ba)|}{\lamsq}
\label{eq:pred}
\end{equation}
reproduces the observed relative miss of every row to within a factor of two,
$18/18$, over three decades of predicted error (Supplementary Material
\S{}S.6, Fig.~S.1 of that section). The five direct misses are the rows where
$\mathrm{d}\lamsq/\mathrm{d}(\ba)$ is largest.

This is a unification, and on its own it would be only an internal one. What
settles the question is an independent third method.

\paragraph{Finite elements settle the disputed cell.}
The most extreme disagreement is at inner~C / outer~F, $\ba=0.840$, $n=0$,
where the table prints $81.0$ and the solver returns a root near $133$. A
criterion was fixed before the finite-element run: a lowest FE mode within
about $2\%$ of $133$ vindicates the solver; one within about $2\%$ of $81$
means the solver misses a real mode.

At C--F, $\ba=0.840$, $n=0$ the lowest flexural FE mode is $\lamsq=132.549$.
The solver gives $132.990$ ($+0.333\%$). Southwell prints $81.0$ ($-63.6\%$).
Nothing in the $40$-mode FE spectrum lies below $132.549$. A refined mesh
(radial divisions $16\to24$, circumferential $256\to384$) reproduces that
value as $132.548973$ against $132.548990$, a relative change of
$1.3\times10^{-7}$. That row is a source error, confirmed externally, and not
a solver defect, a conditioning artifact, a basis truncation, or an
$n$-labelling mistake.

The deck's positive control is the F--C, $\ba=0.5$, $n=0$ cell of Table~2.22,
where the printed value is $17.7$, the solver gives $17.715$ and the finite
elements give $17.7072$: $0.041\%$ against the printed value and $0.044\%$
against the solver. The clamped-edge machinery is therefore correct and the
remaining decks are readable.

\paragraph{Eighteen mixed-edge points.}
Table~\ref{tab:fe18} is the full solver-versus-FE comparison at three C--F
geometries, $n=0$ to $5$, every point sign-ladder CLEAN with $s=0$. All
eighteen agree to better than $0.8\%$; the solver sits \emph{uniformly above}
the finite elements, and the gap grows monotonically with $n$ at every
geometry, from $+0.040\%$ to $+0.757\%$. That is the expected
Kirchhoff-versus-Mindlin signature --- the $4\times4$ is thin-plate while
SHELL281 carries transverse shear, which softens higher-$n$ modes more --- and
it is the same one-sided offset reported in \cite{mccune2026annular}.
Mixed-sign scatter of this magnitude would have been the worrying result;
one-signed monotone growth is not. Two of Southwell's disputed rows are
covered by this set: at $n=3$, $\ba=0.820$ the printed $100$ is $-10.5\%$ from
the FE value while the solver is $+0.451\%$; and the MATCH control at $n=2$,
$\ba=0.522$ has the printed $16.0$ within $0.081\%$ of the FE value.

\begin{table}[htbp]\centering\small
\caption{Solver against independent SHELL281 \cite{ansys} at three inner
clamped / outer free geometries, $\nu=0.30$. Every point is sign-ladder CLEAN
with $s=0$. The offset is one-signed and grows monotonically in $n$:
Kirchhoff versus Mindlin.}
\label{tab:fe18}
\begin{tabular}{lccccccccc}
\toprule
& \multicolumn{3}{c}{$\ba=0.522$} & \multicolumn{3}{c}{$\ba=0.820$}
& \multicolumn{3}{c}{$\ba=0.840$}\\
\cmidrule(lr){2-4}\cmidrule(lr){5-7}\cmidrule(lr){8-10}
$n$ & FE & solver & rel.\ (\%) & FE & solver & rel.\ (\%)
& FE & solver & rel.\ (\%)\\
\midrule
0 & 14.2651 & 14.2708 & $+0.040$ & 104.4191 & 104.6938 & $+0.263$ & 132.5490 & 132.9899 & $+0.333$ \\
1 & 14.5602 & 14.5707 & $+0.072$ & 105.0768 & 105.3775 & $+0.286$ & 133.2228 & 133.6945 & $+0.354$ \\
2 & 15.9871 & 16.0115 & $+0.153$ & 107.0848 & 107.4618 & $+0.352$ & 135.2700 & 135.8325 & $+0.416$ \\
3 & 19.7175 & 19.7624 & $+0.228$ & 110.5422 & 111.0411 & $+0.451$ & 138.7636 & 139.4725 & $+0.511$ \\
4 & 26.5058 & 26.5756 & $+0.263$ & 115.5950 & 116.2545 & $+0.570$ & 143.8146 & 144.7185 & $+0.629$ \\
5 & 36.3693 & 36.4717 & $+0.282$ & 122.4130 & 123.2632 & $+0.695$ & 150.5566 & 151.6958 & $+0.757$ \\
\bottomrule
\end{tabular}
\end{table}

A secondary differential test was also run, comparing the residual radius-ratio
discrepancies of Southwell's table with those of Vogel's on the same footing.
It passed a bar fixed in advance, but its own distribution is heavy-tailed and
it does not control for radius ratio, which the two tables sample differently;
it is reported in \S{}S.6 as corroboration only. The finite-element result,
not that test, is what carries the conclusion of this subsection.

\FloatBarrier
\section{In-plane free--free ring}\label{sec:inplane}
The in-plane reduction of \S\ref{sec:ipring} is compared with Table~2 of
Irie, Yamada and Muramoto \cite{irie1984} at $\nu=0.3$ and
$\beta=0.2,0.4,0.6,0.8$: the first two frequencies at $n=1,2,3,4$ together
with the $n=0$ radial and $n=0$ torsional pairs, $48$ cells in all.
\textbf{All $48$ are MATCH}, the largest relative miss being $0.222\%$ at
$\beta=0.8$, $n=2$, first mode (Table~\ref{tab:irie}; the full $48$-row list
is \S{}S.4). Both $n=0$ families are recovered, which matters because the
radial and torsional branches interleave differently as $\beta$ grows.
A handful of those MATCH cells classify NOT on the sign-flip ladder
($\beta=0.4$ and $0.6$ in Table~\ref{tab:irie}); scoring is from the relative
miss alone, as frozen in \S\ref{sec:screen}, and no in-plane finite-element
check is claimed.

Comparison here is frequency-only. No in-plane eigenvector cross-check is
presented and none is promised: the same comparison was attempted for
rectangular in-plane modes in \cite{mccune2026rect} and was not informative,
for reasons that apply unchanged to this geometry. The in-plane disk is not
claimed: Irie's ``circular'' column is $\beta=0.01$, a tiny-hole control, and
until the in-plane $2\times2$ is built and agrees with it, calling that column
$R_i=0$ would be an assumption rather than a result. The out-of-plane disk of
\S\ref{sec:disk-results} is a genuine $R_i=0$ result, so the omission is
specific to the in-plane motion.

\begin{table}[htbp]\centering\small
\caption{In-plane free--free ring against Irie, Yamada and Muramoto
Table~2 \cite{irie1984}, $\nu=0.3$, with
$\lambda=\omega a\sqrt{\rho(1-\nu^{2})/E}$ of \eqref{eq:irielam}. Twelve
cells per radius ratio; all $48$ MATCH. ``worst'' names the branch carrying
the largest relative miss at that $\beta$.}
\label{tab:irie}
\begin{tabular}{lccccc}
\toprule
$\beta$ & cells & gate & worst branch & worst rel.\ (\%) & ladder classes\\
\midrule
0.2 & 12 & MATCH & $n=1$ & 0.050 & CLEAN,FLOOR \\
0.4 & 12 & MATCH & $n=2$ & 0.053 & CLEAN,FLOOR,NOT \\
0.6 & 12 & MATCH & $n=2$ & 0.031 & CLEAN,FLOOR,NOT \\
0.8 & 12 & MATCH & $n=2$ & 0.222 & CLEAN,FLOOR \\
\bottomrule
\end{tabular}
\end{table}

\FloatBarrier
\section{Discussion}\label{sec:discussion}
What travels from the two companion papers is the Frobenius radial content,
the conversion hygiene of \S\ref{sec:conv}, independent finite elements as a
third method rather than a fitted target, and the discipline of freezing a
screen before the production sweep. What does not travel is the sector
residual bar --- which rejects five genuine full-ring roots --- the sector
subdomain instrument, the rectangular persistence screen, and the Kirchhoff
corner jump, which has no $\theta$-edge to live on.

The limits of the present tables are worth stating plainly.

\begin{itemize}\itemsep3pt
\item Some SP-160 tables state $\nu=1/3$ in their captions and some state no
      Poisson ratio at all. Every sweep that touches an unlabelled table was
      run at both $\nu=0.30$ and $\nu=1/3$, and the two are never mixed in a
      single comparison.
\item Table~2.35 does not print $n=4$; the finite elements carry that branch
      at $\ba=0.5$ as $(4,0)$ and it appears here only through the fidelity
      anchors.
\item The $n=0$ column of Table~2.30 at $\ba=0.9$ is excluded as a source
      defect, on grounds internal to Vogel's own tables. The paper score is
      $72/76$.
\item The $(1,0)$ misses of Table~2.30 at $\ba=0.1$ and $0.3$ are open and
      named.
\item Southwell's Table~2.31 is $13/18$ direct. The sensitivity analysis
      explains the split; the finite elements, not that analysis, are what
      resolve the extreme cell.
\item The disk tiny-hole control at $10^{-3}$ is a named FAIL and its bar is
      not loosened.
\item Search-window endpoint returns are labelled and excluded, never quoted
      as frequencies.
\item Table~2.18 leaves $(2,0)$ and $(2,1)$ at $\ba=0.9$ blank; those cells
      are not invented.
\item Table~2.26 prints $933$ at $(2,1)$, $\ba=0.3$. That is a lost decimal;
      the $s=1$ root is $93.371$. The printed $933$ is a named typeset skip,
      not a MATCH.
\item Table~2.33 misses $(1,0)$ at $\ba=0.1$ at both Poisson ratios
      ($4.72\%$ and $6.02\%$). Named, bar unchanged, not sent to finite
      elements.
\end{itemize}

\section{Conclusions}\label{sec:conclusions}
A validated annular-sector solver reduces, without a new differential
equation, to the classical closed ring and solid disk. The $4\times4$ radial
determinant recovers all nine of Vogel and Skinner's inner/outer combinations
at $704/710$ MATCH${}+{}$WEAK, Narita's polar-orthotropic free--free
Table~3 at $10/10$, and Irie, Yamada and Muramoto's in-plane free--free
table; the $2\times2$ recovers the free solid disk, and the $4\times4$
evaluated at $R_i=0$ demonstrably does not. Independent finite elements
confirm the flexural spine at three free--free geometries, on the disk, and
at three mixed-edge geometries, with a one-sided Kirchhoff-versus-Mindlin
offset below $0.8\%$ throughout --- including the cell where a century-old
table prints $81$ and the plate vibrates at $133$. Piezoelectric
identification, using the free--free and clamped-free ring as its elastic
forward model, is later work.

\small
\section*{Data and code availability}
\begin{sloppypar}
The solver is released under an MIT license as a public GitHub repository,
\url{https://github.com/McCune1/\allowbreak plate_solver}.
For this paper that tree includes the ring and disk entry point
(\texttt{plate\_solver/ring\_disk.py}), the sign-flip ladder and
window-endpoint labelling, the \texttt{TestRingDisk} regression suite that
keeps the disk on the $2\times2$ path, and the tabulated
solver-versus-literature JSON under \texttt{data/paper3/}. Cluster probe
scripts and mixed-edge APDL decks remain in the working package that
produced the tables; they are not part of the Paper~1 or Paper~2 tags.
Tables here were captured under
\texttt{SOLVER\_VERSION} \texttt{2026-07-10.s10}, the same solver as the two
companion papers; results under a different version should be checked against
the clamped-free regression gates first. Printed literature values are
transcribed from the sources cited and are held in one shared module, so that
every table cell in this paper is either a transcription or a computed result
and none is retyped by hand. The Bessel span is Supplementary Material,
\S{}S.1; disk regularity and the rank-3 trap are \S{}S.2; the conversion
algebra is \S{}S.3; the full match lists at the second Poisson ratio and the
remaining Vogel combinations are \S{}S.4; cluster job numbers are \S{}S.5;
the Southwell inverse radius-ratio analysis is \S{}S.6.
\end{sloppypar}

\section*{Declaration of competing interests}
The authors declare that they have no known competing financial interests or
personal relationships that could have appeared to influence the work reported
in this paper.

\section*{Funding}
This research did not receive any specific grant from funding agencies in the
public, commercial, or not-for-profit sectors.

\section*{Acknowledgments}
The computation for this work was performed on The Mill HPC cluster at
Missouri University of Science and Technology \cite{mill2024}. GPU resources
were not used.

\section*{CRediT authorship contribution statement}
\textbf{Garrek McCune:} Conceptualization, Methodology, Software, Validation,
Formal analysis, Investigation, Writing -- original draft, Writing -- review \&
editing, Visualization. \textbf{Daniel Stutts:} Conceptualization,
Supervision.

\section*{Declaration of generative AI and AI-assisted technologies}
During the preparation of this work, the authors used generative AI tools to
assist with manuscript and code organization, drafting, editing, and
formatting. After using these tools, the authors reviewed and edited the
content as needed and take full responsibility for the content of this
publication.
\normalsize

\begin{thebibliography}{99}\small
\setlength{\itemsep}{0pt}
\setlength{\parsep}{0pt}
\bibitem{seok1} J. Seok, H.F. Tiersten, ``Free vibrations of annular sector
cantilever plates. Part 1: out-of-plane motion,'' J. Sound Vib. 271 (2004)
757--772. \url{https://doi.org/10.1016/S0022-460X(03)00414-0}.
\bibitem{seok2} J. Seok, H.F. Tiersten, ``Free vibrations of annular sector
cantilever plates. Part 2: in-plane motion,'' J. Sound Vib. 271 (2004)
773--787. \url{https://doi.org/10.1016/S0022-460X(03)00415-2}.
\bibitem{seok3} J. Seok, H.F. Tiersten, H.A. Scarton, ``Free vibrations of
rectangular cantilever plates. Part 1: out-of-plane motion,'' J. Sound
Vib. 271 (2004) 131--146. \url{https://doi.org/10.1016/S0022-460X(03)00365-1}.
\bibitem{seok4} J. Seok, H.F. Tiersten, H.A. Scarton, ``Free vibrations of
rectangular cantilever plates. Part 2: in-plane motion,'' J. Sound Vib.
271 (2004) 147--158. \url{https://doi.org/10.1016/S0022-460X(03)00366-3}.
\bibitem{leissaSP160} A.W. Leissa, \emph{Vibration of Plates}, NASA SP-160,
Office of Technology Utilization, National Aeronautics and Space
Administration, U.S. Government Printing Office, Washington, D.C., 1969.
\bibitem{vogelSkinner} S.M. Vogel, D.W. Skinner, ``Natural frequencies of
transversely vibrating uniform annular plates,'' J. Appl. Mech. 32 (1965)
926--931. The values compared here are as compiled in \cite{leissaSP160},
Tables~2.22, 2.30 and~2.35; the original paper was not re-read for this work.
\bibitem{southwell} R.V. Southwell, ``On the free transverse vibrations of a
uniform circular disc clamped at its centre; and on the effects of rotation,''
Proc. Roy. Soc. London Ser.~A 101 (1922) 133--153. The values compared here
are as compiled in \cite{leissaSP160}, Table~2.31; the original paper was not
re-read for this work.
\bibitem{raju} P.N. Raju, ``Vibrations of annular plates,'' J. Aeron. Soc.
India 14 (2) (1962) 37--52. Cited via \cite{leissaSP160}, Tables~2.21,
2.29 and~2.34.
\bibitem{narita1984} Y. Narita, ``Natural frequencies of completely free
annular and circular plates having polar orthotropy,'' J. Sound Vib. 92 (1)
(1984) 33--38.
\bibitem{irie1984} T. Irie, G. Yamada, Y. Muramoto, ``Natural frequencies of
in-plane vibration of annular plates'' (Letter to the Editor), J. Sound Vib.
97 (1) (1984) 171--175.
\bibitem{mccune2026annular} G. McCune, D. Stutts, ``Exact wave-function
solution and spurious-mode discrimination for completely free annular sector
plates,'' companion paper, submitted (2026).
\bibitem{mccune2026rect} G. McCune, D. Stutts, ``Exact wave-function solution
of completely free rectangular plates: four-corner Kirchhoff jump and mode
screening,'' companion paper, submitted (2026).
\bibitem{ansys} ANSYS, Inc., \emph{ANSYS Mechanical APDL}, Release 2025 R1,
ANSYS, Inc., Canonsburg, PA, USA (2025). Used for finite-element validation
only.
\bibitem{mill2024} Missouri University of Science and Technology, ``The Mill
HPC Cluster,'' \emph{The Mill} (2024), Article 1.
\url{https://scholarsmine.mst.edu/the-mill/1},
\url{https://doi.org/10.71674/PH64-N397}.
\end{thebibliography}
\end{document}
